\newpage
\section{Aufgaben mit Schwerpunkt Funktionenraum (Kapitel 1)}

\section*{Aufgabe 1: Kurzfragen (12 Punkte)}

\begin{enumerate}[label=\arabic*.]
    \item Die Fourier-Reihe einer $2\pi$-periodischen Funktion $f \in R([-\pi, \pi])$ ist gegeben durch
    \[
        F_\infty(x) = \sum_{k=-\infty}^{\infty} c_k e^{ikx}.
    \]
    Geben Sie die Integralformel für die Koeffizienten $c_k$ an.
    \vspace{0.2cm}
    
    \item Formulieren Sie die Parsevalsche Gleichung (Vollständigkeitsrelation) für die komplexe Fourier-Reihe unter Verwendung der Koeffizienten $c_k$.
    \vspace{0.2cm}
    
    \item Zeigen Sie die Orthogonalitätsrelation für die Exponentialfunktionen:
    \[
        \frac{1}{2\pi} \int_{-\pi}^{\pi} e^{imx} \overline{e^{inx}} \, dx = \delta_{mn}
    \]
    wobei $\delta_{mn}$ das Kronecker-Delta bezeichnet.
    \vspace{0.2cm}
    
    \item Gegeben sei die trigonometrische Fourier-Reihe:
    \[
        F_\infty(x) = \frac{a_0}{2} + \sum_{k=1}^{\infty} \left( a_k \cos(kx) + b_k \sin(kx) \right).
    \]
    Wie hängen die Koeffizienten $c_k$ der komplexen Darstellung mit $a_k$ und $b_k$ zusammen für $k > 0$?
    \vspace{0.2cm}
    
    \item Eine Funktion $f: [-\pi, \pi] \to \mathbb{R}$ sei gerade ($f(-x) = f(x)$). Welche Fourier-Koeffizienten $b_k$ verschwinden? Begründen Sie kurz.
    \vspace{0.2cm}
    
    \item Sei $f(x) = x$ für $x \in (-\pi, \pi)$ und $2\pi$-periodisch fortgesetzt. Ist die Fourier-Reihe von $f$ punktweise konvergent im Punkt $x = \pi$? Wenn ja, welchen Wert nimmt sie dort an?
    \vspace{0.2cm}
    
    \item Definieren Sie den Begriff des \textbf{orthogonalen Funktionensystems} $(\phi_n)_{n \in \mathbb{N}}$ auf einem Intervall $[a, b]$ bezüglich eines Skalarprodukts $\langle \cdot, \cdot \rangle$.
    \vspace{0.2cm}
    
    \item Was besagt der Satz über die \textbf{gleichmäßige Konvergenz} einer Fourier-Reihe? Nennen Sie hinreichende Bedingungen an $f$.
    \vspace{0.2cm}
    
    \item Berechnen Sie den Fourier-Koeffizienten $c_0$ für die Funktion $f(x) = x^2$ auf $[-\pi, \pi]$.
    \vspace{0.2cm}
    
    \item Für welche Werte von $k \in \mathbb{Z}$ gilt $c_k = 0$ für eine reelle, gerade Funktion $f$?
    \vspace{0.2cm}
    
    \item Welche Beziehung besteht zwischen den $L^2$-Norm einer Funktion und ihren Fourier-Koeffizienten laut Parseval-Gleichung?
    \vspace{0.2cm}
    
    \item Sei $f$ eine $2\pi$-periodische Funktion mit Fourier-Koeffizienten $c_k$. Wie verhalten sich die Koeffizienten von $f'(x)$ zu denen von $f(x)$?
    \vspace{0.2cm}
\end{enumerate}

\newpage

\section*{Aufgabe 2: Fourier-Berechnungen (12 Punkte)}
\textbf{(a)} Bestimmen Sie die komplexen Fourier-Koeffizienten $c_k$ der $2\pi$-periodischen Funktion
\[
    f(x) = \begin{cases}
        1 & \text{für } x \in (0, \pi) \\
        -1 & \text{für } x \in (-\pi, 0) \\
        0 & \text{für } x = 0, \pm\pi
    \end{cases}
\]
(Signalfunktion / Signum-Funktion auf $[-\pi, \pi]$)

\vspace{0.2cm}

\textbf{(b)} Nutzen Sie Ihre Ergebnisse aus (a) und die Parsevalsche Gleichung, um die Summe
\[
    \sum_{k=0}^{\infty} \frac{1}{(2k+1)^2}
\]
zu berechnen.

\vspace{0.2cm}

\textbf{(c)} Stellen Sie die Fourier-Reihe Ihrer Funktion aus (a) in der \textbf{trigonometrischen Form} dar und vereinfachen Sie soweit möglich.

\vspace{0.2cm}

\newpage

\section*{Aufgabe 3: Eigenschaften und Konvergenz (12 Punkte)}
\textbf{(a)} Sei $f \in C^1([-\pi, \pi])$ und $2\pi$-periodisch. Zeigen Sie, dass die Fourier-Koeffizienten $c_k$ die Abschätzung
\[
    |c_k| \leq \frac{C}{|k|} \quad \text{für } k \neq 0
\]
erfüllen, wobei $C$ eine Konstante, die von $\|f'\|_\infty$ abhängt.

\vspace{0.2cm}

\textbf{(b)} Betrachten Sie die Funktion $f(x) = |x|$ für $x \in [-\pi, \pi]$ und $2\pi$-periodisch fortgesetzt.

\begin{enumerate}[label=(\roman*)]
    \item Argumentieren Sie, warum alle $b_k = 0$ sind.
    \vspace{0.2cm}
    
    \item Berechnen Sie $a_0$ und $a_k$ für $k \geq 1$.
    \vspace{0.2cm}
    
    \item Geben Sie die Fourier-Reihe explizit an.
    \vspace{0.2cm}
\end{enumerate}

\textbf{(c)} Untersuchen Sie die Konvergenz der Fourier-Reihe von $f(x) = |x|$:
\begin{itemize}
    \item Konvergiert sie punktweise auf $[-\pi, \pi]$?
    \item Konvergiert sie gleichmäßig auf $[-\pi, \pi]$?
\end{itemize}
Begründen Sie Ihre Antworten.

\vspace{0.2cm}

\newpage

\section*{Aufgabe 4: Orthogonalität und Projektionen (12 Punkte)}
\textbf{(a)} Zeigen Sie, dass das Funktionensystem $\{1, \cos(x), \sin(x), \cos(2x), \sin(2x), \ldots\}$ orthogonal bezüglich des Skalarprodukts
\[
    \langle f, g \rangle = \int_{-\pi}^{\pi} f(x)g(x) \, dx
\]
auf dem Raum der stückweise stetigen Funktionen ist. Berechnen Sie jeweils $\|1\|^2$, $\|\cos(kx)\|^2$, und $\|\sin(kx)\|^2$.

\vspace{0.2cm}

\textbf{(b)} Sei $V_n = \operatorname{span}\{1, \cos(x), \sin(x), \ldots, \cos(nx), \sin(nx)\}$. Sei $f \in L^2([-\pi, \pi])$. Zeigen Sie, dass die Partialsumme $S_n f$ der Fourier-Reihe die \textbf{bestmögliche Approximation} von $f$ in $V_n$ bezüglich der $L^2$-Norm ist, d.h.:
\[
    \|f - S_n f\|_2 \leq \|f - p\|_2 \quad \forall p \in V_n
\]

\vspace{0.2cm}

\textbf{(c)} Gegeben sei $f(x) = x^2$ auf $[-\pi, \pi]$. Berechnen Sie die quadratische Approximationsgüte
\[
    \|f - S_1 f\|_2^2 = \int_{-\pi}^{\pi} |f(x) - S_1 f(x)|^2 \, dx
\]
wobei $S_1 f$ die erste Partialsumme ist.

\vspace{0.2cm}

\newpage

\section*{Aufgabe 5: Anwendungen und Vertiefung (12 Punkte)}
\textbf{(a)} Die Funktion $f(x) = e^x$ sei auf $[-\pi, \pi]$ betrachtet und $2\pi$-periodisch fortgesetzt. Berechnen Sie die komplexen Fourier-Koeffizienten $c_k$.

\vspace{0.2cm}

\textbf{(b)} Verwenden Sie die Parsevalsche Gleichung auf die Funktion aus (a), um die Identität
\[
    \sum_{k=-\infty}^{\infty} \frac{1}{1+k^2} = \text{???}
\]
zu überprüfen bzw. den korrekten Wert zu bestimmen.

\vspace{0.2cm}

\textbf{(c)} Sei $(f_n)_{n \in \mathbb{N}}$ eine Folge von Funktionen in $L^2([-\pi, \pi])$, die \textbf{orthonormal} ist. Zeigen Sie die \textbf{Besselsche Ungleichung}:
\[
    \sum_{n=1}^{\infty} |\langle f, f_n \rangle|^2 \leq \|f\|_2^2
\]
Für welche $f$ wird hier Gleichheit erreicht?